
% ==========================================
% PHAN C - OSPF TRONG AS B
% ==========================================

\section{Phần C -- Cấu hình OSPF trong AS B}

% ==========================================
% A. MUC TIEU
% ==========================================

\subsection{Mục tiêu}

Phần này triển khai giao thức định tuyến
OSPF (Open Shortest Path First) trong
hệ thống AS B, bao gồm ba router
B1, B2 và B3.

\textbf{Các yêu cầu kỹ thuật cần đạt:}

\begin{itemize}
    \item Cấu hình OSPF process 1,
    area 0 trên cả ba router B1, B2 và B3.

    \item Thiết lập OSPF trên ba liên kết
    nội bộ L6, L7 và L8 để các router
    có thể hình thành quan hệ láng giềng
    và trao đổi thông tin định tuyến.

    \item Kiểm chứng trạng thái
    OSPF Neighbor và các tuyến
    học được thông qua OSPF.
\end{itemize}


% ==========================================
% B. TRIEN KHAI CAU HINH
% ==========================================

\subsection{Triển khai cấu hình}

\subsubsection{Nguyên tắc cấu hình OSPF}

AS B sử dụng giao thức OSPF làm giao thức
định tuyến nội bộ (IGP), bao gồm ba router
B1, B2 và B3 kết nối theo topology tam giác
thông qua các liên kết L6, L7 và L8.

Toàn bộ các liên kết nội bộ được cấu hình
trong OSPF area 0. Mỗi router sử dụng
địa chỉ Loopback0 làm Router ID.

Nhóm sử dụng các câu lệnh \texttt{network}
để chỉ định những mạng tham gia OSPF.
Cụ thể, OSPF được kích hoạt trên các
liên kết nội bộ L6--L8, các Loopback0
và mạng LAN-B.

\subsubsection{Giải thích các lệnh cấu hình OSPF chính}

Trong quá trình triển khai OSPF trên các router
B1, B2 và B3, nhóm sử dụng các lệnh cấu hình
chính như sau:

\begin{itemize}

    \item \textbf{\texttt{router ospf 1}:}
    Khởi tạo tiến trình định tuyến OSPF với
    Process ID là 1. Process ID chỉ có ý nghĩa
    cục bộ trên từng router, không bắt buộc
    giống nhau giữa các router để hình thành
    quan hệ láng giềng.

    \item \textbf{\texttt{router-id A.B.C.D}:}
    Thiết lập Router ID để định danh duy nhất
    một router trong hệ thống OSPF.
    Trong đồ án, Router ID được đặt bằng
    địa chỉ Loopback0 của từng router nhằm
    đảm bảo tính ổn định.

    \item \textbf{\texttt{network A.B.C.D W.X.Y.Z area 0}:}
    Chỉ định các interface có địa chỉ IP
    phù hợp với địa chỉ mạng và wildcard mask
    để tham gia tiến trình OSPF trong area 0.

    Wildcard mask được xác định bằng cách
    lấy 255.255.255.255 trừ subnet mask.

    Ví dụ, với mạng 10.71.26.20/30,
    subnet mask là 255.255.255.252,
    do đó wildcard mask là 0.0.0.3.

    \begin{lstlisting}
network 10.71.26.20 0.0.0.3 area 0
    \end{lstlisting}

    Lệnh trên kích hoạt OSPF trên interface
    có địa chỉ thuộc mạng 10.71.26.20/30,
    tương ứng với liên kết L6 giữa B1 và B2.

    \item \textbf{\texttt{ip ospf network point-to-point}:}
    Thiết lập kiểu mạng OSPF của interface
    thành point-to-point.

    Trong đồ án, lệnh này được sử dụng
    trên Loopback1 của B2 nhằm đảm bảo
    mạng LAN-B được quảng bá với prefix
    /26 thay vì host route /32 như
    hành vi mặc định của OSPF trên Loopback.

    \begin{lstlisting}
interface Loopback1
 ip ospf network point-to-point
    \end{lstlisting}

\end{itemize}


% ==========================================
% CAU HINH B1
% ==========================================

\subsubsection{Cấu hình OSPF trên router B1}

\begin{lstlisting}
config t
router ospf 1
 router-id 10.171.26.4
 network 10.71.26.20 0.0.0.3 area 0
 network 10.71.26.28 0.0.0.3 area 0
 network 10.171.26.4 0.0.0.0 area 0
end
\end{lstlisting}

% ==========================================
% CAU HINH B2
% ==========================================

\subsubsection{Cấu hình OSPF trên router B2}


\begin{lstlisting}
config t
interface Loopback1
 ip ospf network point-to-point
exit

router ospf 1
 router-id 10.171.26.5
 network 10.71.26.20 0.0.0.3 area 0
 network 10.71.26.24 0.0.0.3 area 0
 network 10.171.26.5 0.0.0.0 area 0
 network 10.171.126.64 0.0.0.63 area 0
end
\end{lstlisting}

% ==========================================
% CAU HINH B3
% ==========================================

\subsubsection{Cấu hình OSPF trên router B3}


\begin{lstlisting}
config t
router ospf 1
 router-id 10.171.26.6
 network 10.71.26.24 0.0.0.3 area 0
 network 10.71.26.28 0.0.0.3 area 0
 network 10.171.26.6 0.0.0.0 area 0
end
\end{lstlisting}



% ==========================================
% C. KIEM CHUNG KET QUA
% ==========================================

\subsection{Kiểm chứng kết quả}

\subsubsection{Kiểm tra quan hệ OSPF Neighbor}

\begin{table}[H]
\centering
\caption{Các quan hệ OSPF Neighbor cần kiểm chứng}
\label{tab:partC-neighbor}
\small
\begin{tabularx}{\textwidth}{c c X c}
\toprule
\textbf{Router} &
\textbf{Neighbor} &
\textbf{Liên kết} &
\textbf{Trạng thái yêu cầu} \\
\midrule
B1 & B2 & L6 & FULL \\
B1 & B3 & L8 & FULL \\
B2 & B1 & L6 & FULL \\
B2 & B3 & L7 & FULL \\
B3 & B1 & L8 & FULL \\
B3 & B2 & L7 & FULL \\
\bottomrule
\end{tabularx}
\end{table}

Trạng thái FULL cho biết
hai router đã hoàn tất quá trình
đồng bộ cơ sở dữ liệu trạng thái
liên kết (LSDB).

Mỗi quan hệ láng giềng được
hiển thị từ cả hai phía, do đó
ba liên kết nội bộ tương ứng
với ba cặp quan hệ Neighbor.

% ------------------------------------------
% HINH OSPF NEIGHBOR
% ------------------------------------------

\partzeroimage
{images/configuration/part_C/B1_neighbor.png}
{Danh sách OSPF Neighbor trên router B1}
{fig:partC-B1-neighbor}

\partzeroimage
{images/configuration/part_C/B2_neighbor.png}
{Danh sách OSPF Neighbor trên router B2}
{fig:partC-B2-neighbor}

\partzeroimage
{images/configuration/part_C/B3_neighbor.png}
{Danh sách OSPF Neighbor trên router B3}
{fig:partC-B3-neighbor}

\newpage

% ==========================================
% KIEM TRA ROUTING TABLE
% ==========================================

\subsubsection{Kiểm tra bảng định tuyến OSPF}

Để kiểm tra các tuyến đường
được học thông qua OSPF,
thực hiện lệnh sau trên B2:

\begin{lstlisting}
show ip route ospf
\end{lstlisting}

Các tuyến học được thông qua
OSPF sẽ có ký hiệu \texttt{O}
trong bảng định tuyến.

Đối với B2, những mạng đã kết nối
trực tiếp qua L6 và L7 được hiển thị
dưới dạng Connected (C),
không phải tuyến OSPF.

Ngược lại, mạng L8 kết nối
giữa B1 và B3 và các Loopback0
của B1, B3 có thể được học
thông qua OSPF.

Các thông tin quan trọng
trong bảng định tuyến gồm:

\begin{itemize}
    \item \textbf{Destination:}
    Mạng đích được học.

    \item \textbf{Administrative Distance:}
    Giá trị AD mặc định của OSPF
    là 110.

    \item \textbf{Metric:}
    Chi phí OSPF để đến mạng đích.

    \item \textbf{Next-hop:}
    Địa chỉ router kế tiếp
    trên đường đến mạng đích.
\end{itemize}

\partzeroimage
{images/configuration/part_C/B2_route_ospf.png}
{Các tuyến đường OSPF học được trên router B2}
{fig:partC-B2-route}

\newpage